\section{Discussion and conclusion}
\label{sec:conclusion}

The main message of this paper is simple: in repeated-measurement QRC, \emph{shot organization is part of the learning algorithm}. Split-ensemble training sits between the two standard extremes. Relative to EV, it increases the number of readout-training rows. Relative to Raw, it reduces the feature noise seen by the readout. The exact grouped-view lift $\rhoK=(\Nshots/k)\rhoEV$ makes this tradeoff transparent, and the attenuation-versus-row-count surrogate explains why an intermediate $k$ is often preferred.

The empirical evidence now has four layers of support. First, the empirical $\rhoEV$ plot shows the regime picture directly: the gain is largest when EV is still row-limited and decays as EV becomes better conditioned. Second, the baseline and duplication controls show that the effect is not explained away by simple row-count matching or by the diagonal noise-aware proxy. Third, the larger-reservoir and architecture sweeps show that the gain is not tied to a single tiny ring-CNOT reservoir. Fourth, the same qualitative ranking survives on real hardware.

A practical decision rule follows immediately.
\begin{enumerate}[leftmargin=*]
    \item Compute $\rhoEV=\Ttrain/d$ under exact budget accounting, including the factor $\Rsettings$.
    \item If $\rhoEV$ is comfortably above the conditioning-limited regime, EV is often already sufficient.
    \item If $\rhoEV$ is small or moderate, validate over the divisors of $\Nshots$, using Eq.~\eqref{eq:kstar} as a warm start for the preferred $k$.
    \item Keep inference aggregated at the labeled time-step level; the protocol changes training, not the basic prediction interface.
\end{enumerate}

The paper also has clear limits. We do not claim that split-ensemble universally dominates all estimator families. The noise-aware baselines in the present study are diagonal surrogates rather than a full implementation of Ahmed, Tennie, and Magri \cite{ahmed2025optimal}. We also do not claim that the same numeric $\rhoEV$ crossover transfers unchanged across all reservoirs, readouts, or datasets. Finally, we have stayed with linear ridge readouts to isolate the conditioning effect; a next step is to test how much of the protocol advantage carries over to richer readout families.

Within that scope, however, the result is strong. The method is budget fair, easy to implement, compatible with standard inference, theoretically interpretable, validated on hardware, and robust across both reservoir size and reservoir architecture. For a QCE audience, that is the key point: if shots are expensive, then how we \emph{organize} them is not an afterthought but a real algorithmic choice.
